\section*{Zusammenfassung}

Lärm- und Schwingungsminderung an leichten Strukturen ist eine verbreitete ingenieurtechnische Aufgabe. Platten in Autos, Flugzeugen und Maschinen brauchen oft zusätzliche Masse oder Dämpfung, um die abgestrahlte Schallleistung zu verringern, und eine leichte Möglichkeit dafür ist, wenige kleine, abgestimmte Masse-Resonatoren an der Platte zu befestigen. Jeder Resonator ist eine kleine Kombination aus Feder, Dämpfer und Masse an einer Position auf der Platte, und zusammen bestimmen diese Resonatoren, wie viel Schallleistung die Platte abstrahlt. Diese Arbeit beschreibt die abgestrahlte Schallleistung mit der äquivalenten Strahlungsleistung (ERP), ausgewertet zwischen 10 und 160 Hz.

Eine gute Anordnung der Resonatoren von Hand zu finden ist schwierig, weil die Anzahl der Resonatoren, ihre Werte und ihre Positionen alle zusammenwirken. Eine exakte, geschlossene Lösung existiert nur für wenige einfache Fälle; die meisten realistischen Szenarien benötigen eine vollständige Mehrfeldsimulation, die zwar genau, aber viel zu langsam ist, um tausende Anordnungen durchzuprobieren, und die nur eine Richtung des Problems beantwortet: von einem Entwurf zu seinem Spektrum. Was ein Ingenieur eigentlich braucht, ist auch die andere Richtung: ausgehend von einem gewünschten akustischen Verhalten einen passenden Resonator-Entwurf vorzuschlagen. Das nennt man generatives, oder inverses, Design, und es ist das Hauptziel dieser Arbeit.

Um dieses Ziel zu erreichen, wird zunächst mit einem validierten Platten-Löser ein Datensatz aus Resonator-Entwürfen und ihren ERP-Spektren erzeugt. Jeder Resonator wird durch fünf Größen beschrieben, Masse, Abstimmfrequenz, Position und ein Dämpfungsverhältnis, die mit Latin-Hypercube-Sampling gezogen werden, damit der Datensatz den Entwurfsraum gleichmäßig abdeckt. Zehn neuronale Netzwerk-Architekturen werden als schnelle Vorwärts-Surrogate auf einer Zwei-Resonator-Version dieses Datensatzes trainiert und sagen das ERP-Spektrum eines Entwurfs in einem Bruchteil der Zeit einer vollständigen Simulation voraus; das beste Modell erreicht einen mittleren quadratischen Fehler von etwa 1,1\,dB auf ungesehenen Entwürfen. Dieselben zehn Architekturen werden außerdem auf einem Drei-Resonator-Datensatz verglichen, auf dem sich der Fehler je nach Architektur um das 1,1- bis 2-Fache erhöht (beim besten Modell von etwa 1,1\,dB auf 2,2\,dB), was zeigt, wie viel schwieriger das Problem mit mehr Resonatoren wird. Darauf aufbauend werden invertierbare neuronale Operatoren entwickelt, die mit einem einzigen Satz von Gewichten in beide Richtungen arbeiten: dasselbe Modell sagt ein Spektrum aus einem Entwurf voraus und schlägt umgekehrt einen Entwurf für ein Zielspektrum vor. Diese werden auf einem früheren, größeren Zwei-Resonator-Datensatz ausgewertet und können bereits mehrere Entwurfsgrößen recht gut rekonstruieren, auch wenn die genauen Resonatorpositionen schwer zu rekonstruieren bleiben, da mehrere unterschiedliche Anordnungen fast dasselbe Spektrum erzeugen können. Ein weiteres Experiment prüft außerdem, wie gut die Vorwärtsmodelle auf Frequenzen extrapolieren, die sie beim Training nie gesehen haben, was sich als deutlich schwieriger erweist als die Interpolation innerhalb des Trainingsbereichs.

\newpage
